\section{Related Work}
\label{sec:related_work}

Our work sits at the intersection of several active research areas: methods for building interpretability into models during training, techniques for steering model behavior through representation manipulation, and approaches for removing specific model capabilities.

\subsection{Concept-Based Interpretability Methods}

Several families of methods aim to make neural networks interpretable through human-meaningful concepts. Concept Bottleneck Models~\citep{KohNTMPKL20} enforce interpretability architecturally by introducing an intermediate layer where each dimension corresponds to a predefined concept, enabling test-time interventions---though originally requiring full supervision, recent work has reduced this burden through post-hoc discovery or sparse training-time methods with minimal labels~\citep{oikarinen2023labelfreeconceptbottleneckmodels,semenov2024sparseconceptbottleneckmodels,Sawada2022ConceptBM}. However, post-hoc interpretability methods have revealed that concepts in CBMs may not correspond to semantically meaningful input features~\citep{Margeloiu2021DoCB}, motivating approaches that explicitly structure concept representations during training. Concept Activation Vectors~\citep{KimWGCWVS18} take a lightweight post-hoc approach, learning linear probes from as few as 30 examples per concept to identify where concepts appear in trained models---useful for bias detection but providing no architectural guarantees for interventions. Sparse Autoencoders use unsupervised dictionary learning to discover interpretable features models actually use, recently scaling to frontier language models~\citep{HubenCRES24,templeton2024scaling}, though features are discovered rather than positioned during training. Concept Whitening~\citep{Chen_2020} replaces batch normalization with transformations that align latent space axes with concepts using representative examples, enabling layer-wise interpretability without hurting performance. These methods trade off supervision requirements, timing of concept incorporation (training vs. post-hoc), and intervention capabilities.

\subsection{Machine Unlearning and Representation Engineering}

Machine unlearning addresses removing specific capabilities from trained models, driven by privacy regulations and safety concerns. Gradient-based methods attempt to reverse training through gradient ascent on "forget" data~\citep{jang2022knowledgeunlearningmitigatingprivacy,yao-etal-2024-machine,zhang-etal-2024-negative}, but face challenges with gradient explosion, catastrophic forgetting, and instability---particularly at high forget rates. Representation-based methods like RMU~\citep{li2024wmdpbenchmarkmeasuringreducing} redirect activations of unwanted content toward random directions, reducing hazardous knowledge to near-random performance on benchmarks like WMDP, though the distinction between masking and true removal remains unclear. Training-time approaches remain rare: SISA~\citep{bourtoule-etal-2021-machine} enables efficient removal through data sharding but with substantial computational overhead, while Ready2Unlearn~\citep{duan2025ready2unlearnlearningtimeapproachpreparing} uses meta-learning to prepare models for later unlearning---yet both operate through data organization or optimization dynamics rather than explicit geometric positioning. Representation engineering methods manipulate behavior by modifying internal activations~\citep{zou2025representationengineeringtopdownapproach}: activation addition~\citep{TurnerTUDLMM23} extracts steering vectors from paired prompts, while optimized methods like BiPO~\citep{cao2024personalizedsteeringlargelanguage} and reinforcement learning approaches~\citep{sinii2025smallvectorsbigeffects} train better steering vectors---but all depend on directions discovered in already-trained models. Systematic analysis reveals substantial reliability issues: steering effectiveness varies dramatically across inputs, many concepts prove "anti-steerable", and success often depends on spurious correlations rather than coherent concepts~\citep{TanCLPKGK24}. Abliteration~\citep{arditi2024refusallanguagemodelsmediated} demonstrated that safety behaviors can be removed through targeted weight orthogonalization with negligible performance degradation, providing evidence for the linear representation hypothesis---yet achieving selective ablation without side effects remains challenging when features are distributed or when networks exhibit "compensatory masquerade" by routing around ablations~\citep{meyes-etal-2019-ablation}.

\subsection{Geometric Constraints in Neural Networks}

Recent work has explored how geometric constraints on representations improve training dynamics and enable interpretability. nGPT~\citep{loshchilov2025ngptnormalizedtransformerrepresentation} normalizes all transformer components to unit norm, constraining representations to a hypersphere, yielding 4-20× faster convergence, more interpretable angular relationships, and stable gradients---suggesting hypersphere constraints improve both interpretability and optimization itself. Angular margin losses from face recognition~\citep{liu2018spherefacedeephypersphereembedding,Deng_2022} enforce separation between classes in hyperspherical geometry through L2-normalized features and additive margins, achieving state-of-the-art results because angular constraints create geometrically clean separation. Theoretical analysis shows contrastive learning on hyperspheres naturally optimizes for alignment and uniformity~\citep{wang2022understandingcontrastiverepresentationlearning}---properties that facilitate linear separability and robust representations. In the context of concept-based models, concept orthogonal loss~\citep{Sheth2023AuxiliaryLF} encourages separation between learned concept representations while reducing intra-concept distance, improving concept disentanglement in CBMs through auxiliary training objectives---though applied to dense concept bottlenecks rather than sparse, pre-positioned concepts. While geometric constraints have improved training efficiency and discriminability, their use specifically for positioning concepts to enable interventions---particularly with minimal supervision---remains less explored.

\subsection{Positioning Our Work}

Our method integrates training-time geometric positioning of concepts through hypersphere constraints with sparse supervision ($<0.1\%$ of examples labeled per concept). We build most directly on sparse CBM work that reduces supervision requirements (Sparse-CBM, UCBM, Z-CBM) and on geometric representation learning (nGPT, angular margin methods) that uses hypersphere constraints for improved training and interpretability. The closest related work differs in key trade-offs: CBMs and variants provide training-time structure but vary in supervision needs and intervention capabilities; Concept Whitening uses whitening transformations rather than hypersphere normalization and requires moderate supervision through concept datasets; representation engineering methods discover steering directions post-hoc, inheriting reliability issues from emergent geometry; training-time unlearning methods (Ready2Unlearn, SISA) use meta-learning or data organization rather than concept positioning.

Our observation-driven approach to hyperparameter tuning---watching latent geometry evolve during training to manually adjust time-varying regularizer schedules---represents a form of developmental interpretability. This differs from recent work studying how data distribution changes during training affect final models; we instead manipulate optimization dynamics through coordinated loss term schedules while training on fixed data. The motivation for training-time intervention comes from evidence that neural networks lose plasticity early in training and establish connectivity patterns during critical periods~\citep{AchilleRS17}, making post-hoc modification difficult.

Our method addresses three specific gaps. First, unlike sparse CBMs that discover concepts during training, we \textit{fix} concept locations a priori, enabling pre-planned interventions. Second, unlike CAVs and steering vectors that depend on emergent geometry, we \textit{construct} geometry that supports interventions by design. Third, unlike machine unlearning methods that attempt post-hoc removal, we achieve selective ablation by isolating concepts during training---preventing the entanglement that makes clean removal difficult in standard architectures. Whether this approach scales beyond color autoencoders to complex domains like language models---particularly whether geometric constraints remain tractable in high-dimensional spaces with attention and residual pathways---remains an important open question for future work.
